% =====================================================================
%  ПРИМЕР. Базовое резюме, семейство ролей: backend, рынок: rf, язык: ru
%  Собрано скиллом build-resume из example/facts.md.
%  Человек и компании вымышлены.
% =====================================================================
\documentclass[a4paper,10pt]{article}

\usepackage{cmap}
\usepackage[T2A]{fontenc}
\usepackage[utf8]{inputenc}
\usepackage[russian,english]{babel}
\usepackage{geometry}
\usepackage{enumitem}
\usepackage{hyperref}
\usepackage{xurl}
\usepackage{microtype}
\usepackage{titlesec}

\geometry{left=1.35cm,right=1.35cm,top=1.25cm,bottom=1.25cm}
\setlength{\parindent}{0pt}
\setlist[itemize]{leftmargin=1.15em, itemsep=2pt, topsep=2pt, parsep=0pt}

\titleformat{\section}{\large\bfseries}{}{0em}{}[\titlerule]
\titlespacing{\section}{0pt}{6pt}{4pt}

\hypersetup{colorlinks=true, urlcolor=black, linkcolor=black}
\Urlmuskip=0mu plus 1mu
\emergencystretch=2em
\sloppy
\input{glyphtounicode}
\pdfgentounicode=1

% \mbox защищает от переноса после «/»: иначе извлекается битый токен «CI/ CD».
\newcommand{\CICD}{\mbox{CI/CD}}
\newcommand{\RESTAPI}{\mbox{REST API}}

% ---------- ДАННЫЕ (profile.json, комплект CONTACT-RF) ----------
\newcommand{\CVName}{Алексей Морозов}
\newcommand{\CVTitle}{Backend-разработчик (Python / Go)}
\newcommand{\CVEmail}{a.morozov@example.com}
\newcommand{\CVPhone}{+7 900 000 11 22}
\newcommand{\CVTelegram}{amorozovdemo}

\begin{document}

% FACT: IDENT-name CONTACT-RF TITLE-backend IDENT-telegram LOC-ekb
{\LARGE \textbf{\CVName}}\\
\textbf{\CVTitle}\\
Email: \href{mailto:\CVEmail}{\CVEmail} \quad
Телефон: \CVPhone \quad
Telegram: \href{https://t.me/\CVTelegram}{@\CVTelegram}\\
Екатеринбург \quad Удалённо или гибрид \quad Английский --- B2

\section*{О себе}

% FACT: SUM-backend
Backend-разработчик с 5 годами опыта, из них 3 года в финтехе. Пишу на Python
(FastAPI, asyncio) и Go, проектирую и сопровождаю микросервисы под нагрузкой:
платёжный шлюз, над которым работает моя команда, обрабатывает порядка 400 запросов
в секунду в пике. Отдельная сильная сторона --- PostgreSQL глубже CRUD: разбор планов
запросов, индексы, партиционирование. Довожу задачи до эксплуатации, а не до мержа:
пишу тесты, настраиваю метрики и разбираю инциденты по своим сервисам.

\section*{Навыки}

% FACT: SKL-lang SKL-frameworks SKL-db SKL-infra SKL-arch SKL-obs SKL-test
\textbf{Языки:} Python (5 лет, asyncio, typing), Go (2 года, сервисы в проде), SQL,
Bash, Git; Rust --- базово\\
\textbf{Фреймворки и протоколы:} FastAPI, aiohttp, Celery, gRPC, \RESTAPI{};
Django --- работал 1,5 года\\
\textbf{Базы данных и очереди:} PostgreSQL (планы запросов, индексы,
партиционирование), Redis (кеш, распределённые блокировки); Kafka --- чтение и
запись из рабочих задач; MongoDB --- базово\\
\textbf{Инфраструктура:} Docker, GitLab \CICD{}, Linux; Kubernetes --- выкатка по
готовым манифестам, разбор падений подов\\
\textbf{Архитектура и надёжность:} микросервисная декомпозиция, идемпотентность,
ретраи, обработка частичных отказов; saga --- ограниченно\\
\textbf{Наблюдаемость и тесты:} Prometheus, Grafana, Sentry, структурные логи,
pytest, testcontainers

\section*{Опыт}

% FACT: EXP-aurora
\textbf{Аврора Финтех} \hfill \textit{март 2023 --- н.в.}\\
\textit{Backend-разработчик}
\begin{itemize}
    \item Разрабатываю и сопровождаю платёжный шлюз на Python (FastAPI, asyncio):
          приём платежей, сверка с провайдерами, повторы и идемпотентность операций.
    \item Спроектировал схему идемпотентности платёжных операций: повторная отправка
          запроса не приводит к двойному списанию.
    \item Перевёл два внутренних сервиса с Python на Go; время ответа на тяжёлых
          ручках снизилось примерно вдвое.
    \item Оптимизировал запросы к PostgreSQL: разбор планов, недостающие индексы,
          партиционирование таблицы транзакций по месяцам.
    \item Настроил метрики и алерты (Prometheus, Grafana) по своим сервисам; дежурю
          в ротации, разбираю инциденты и пишу постмортемы.
    \item Пишу тесты: pytest, testcontainers для интеграционных, покрытие критичных путей.
    \item Провожу код-ревью, менторю стажёров: две стажировки, обе закончились наймом.
\end{itemize}

% FACT: EXP-cityline
\textbf{Ситилайн Ритейл} \hfill \textit{июль 2021 --- февраль 2023}\\
\textit{Backend-разработчик}
\begin{itemize}
    \item Разрабатывал сервисы каталога и корзины интернет-магазина на Python
          (Django, затем FastAPI); \RESTAPI{} для веб- и мобильных клиентов.
    \item Вынес выгрузку прайс-листов в асинхронные задачи (Celery, Redis): выгрузка
          перестала блокировать основной поток и укладывалась в ночное окно.
    \item Написал слой кеширования каталога на Redis; нагрузка на основную базу
          в часы пик заметно снизилась.
    \item Участвовал в распиливании монолита: выделил каталог в отдельный сервис
          вместе с миграцией данных.
    \item Настраивал пайплайны GitLab \CICD{}: прогон тестов, линтеры, сборка образов.
\end{itemize}

% FACT: EXP-cityline-intern
\textbf{Ситилайн Ритейл} \hfill \textit{февраль 2021 --- июнь 2021}\\
\textit{Стажёр-разработчик}
\begin{itemize}
    \item Правил баги в существующем коде, писал юнит-тесты, разбирался с чужой
          кодовой базой. По итогам стажировки получил оффер.
\end{itemize}

\section*{Проекты}

% FACT: PRJ-ratelimit
\textbf{Rate limiter на Go} \hfill \textit{личный проект, открытый код}
\begin{itemize}
    \item Библиотека ограничения частоты запросов: token bucket и sliding window,
          состояние в Redis для распределённого режима.
    \item Бенчмарки и тесты на гонки. Используется в одном из внутренних сервисов
          на работе.
\end{itemize}

% FACT: PRJ-migrator
\textbf{Инструмент миграций схемы БД} \hfill \textit{рабочий, внутренняя библиотека}
\begin{itemize}
    \item CLI поверх Alembic: прогон миграций в нужном порядке для нескольких
          сервисов, откат, проверка на несовместимые изменения перед выкаткой.
    \item Вынесен во внутреннюю библиотеку, используют три команды.
\end{itemize}

\section*{Образование}

% FACT: EDU-utu
\textbf{Уральский технический университет} \hfill \textit{2017 --- 2021}\\
Программная инженерия, бакалавр

% FACT: EDU-go EDU-pg
\textbf{Доп. обучение}
\begin{itemize}
    \item Go для backend-разработчиков --- по итогам перевёл два внутренних сервиса
          с Python на Go.
    \item PostgreSQL: планы запросов и оптимизация --- разбор slow query log,
          переписывание запросов, подбор индексов.
\end{itemize}

\section*{Достижения}

% FACT: AWD-hackathon
\begin{itemize}
    \item Финтех-хакатон 2022 --- 2 место в командном зачёте.
\end{itemize}

\end{document}
